\section*{Introduction}\label{sec:intro-d6}
\addcontentsline{toc}{section}{Introduction}

% D6 brief: Context, background, Company description, member bios/roles, member contributions.
% This section is EXCLUDED from the 15-page limit.

\subsection*{Context and Background}

In 2023, UK mountain rescue teams responded to over 3{,}000 incidents, with search phases frequently lasting several hours across terrain that is difficult or dangerous to traverse on foot~\cite{mra2023}. The \textit{golden hour}---the window within which medical intervention most significantly improves survival---places severe pressure on response teams to locate casualties rapidly~\cite{goodrich2008}. A ground search team of ten can sweep approximately $0.1$--$0.5$\,km$^{2}$/hr; a single UAV at \SI{35}{\metre} altitude with a \ang{60} field of view can survey a 32\,m-wide swath at \SI{8}{\metre\per\second}, covering approximately $0.6$\,km$^{2}$/hr---an order-of-magnitude improvement per searcher deployed~\cite{goodrich2008, murphy2014}. Recent advances in lightweight object-detection models, particularly the YOLO family~\cite{yolohistory, yolov8}, have made it feasible to run real-time inference on edge hardware such as the Raspberry Pi, enabling fully onboard detection without continuous video downlinks~\cite{edgeai}.

This report presents the design, implementation, and evaluation of an autonomous search and rescue (SAR) drone that addresses the gap between expensive commercial platforms (\pounds5\,000--\pounds30\,000 requiring dedicated ground operators) and research prototypes relying on GPU-class computers too heavy for volunteer rescue organisations. The central question is whether a \textit{low-cost, fully autonomous} SAR UAV---built from a Raspberry Pi, an RGB camera, and commodity components---can detect, localise, and respond to a casualty with all decision-making performed onboard.

The project scenario, set by the AENGM0074 module, is as follows. A hiker has gone missing in Fenswood Wilderness, a semi-rural landscape managed by the University of Bristol. Each student company receives identical hardware---a Hexsoon EDU-450 hexacopter, a Cube Orange+ flight controller, a Raspberry Pi~5, and a global-shutter camera---and must build a UAV capable of:

\begin{itemize}
    \item systematically searching a defined area while respecting a no-fly zone over a Site of Special Scientific Interest (SSSI);
    \item detecting a life-size casualty surrogate (``rescue mannequin'') using onboard computer vision;
    \item reporting the casualty's estimated GPS coordinates and imagery to a ground operator;
    \item landing between 5 and 10\,m from the casualty and deploying a simulated first-aid kit; and
    \item returning to the designated take-off location.
\end{itemize}

A 50\,m altitude ceiling, defined flight-area boundaries, and automatic geofencing to enforce SSSI exclusion must all be respected without pilot intervention. The scenario is deliberately constrained to a single outdoor demonstration flight preceded by progressive bench and flight testing, reflecting the compressed timeline of a single-term MSc module. Twelve requirements (R01--R12) govern the mission. All software is released publicly under the MIT licence.

\subsection*{Company Description}

The company comprises five MSc students from the Department of Aerospace Engineering, formed at the start of the AENGM0074 module. We adopted a workstream-based methodology: five parallel streams---Computer Vision, Hardware Integration, Search Logic, Ground Station, and Safety/Testing---fed into defined integration milestones where subsystems were combined and verified before advancing to the next stage. All code resides in a shared GitHub repository\footnote{Repository: \url{https://github.com/DimaChup/Group_Proj} (MIT licence, R12).}, with a single protected main branch and feature branches for individual work.

Day-to-day coordination relied on a group chat for informal updates and weekly in-person meetings for design decisions, integration planning, and blocker resolution. A shared status tracker documented every group decision, assigned action items with owners and deadlines, and maintained a running log of blockers and dependencies. This lightweight process allowed each member to work independently on their stream while keeping integration points well-defined and preventing subsystems from developing in isolation.

\subsection*{Team Members and Roles}

\textbf{Dmytro Chuprynyuk --- Lead Software Developer.}
MSc student in Robotics with a background in software engineering and embedded systems. Designed and implemented the autonomous software stack: finite state machine, computer vision pipeline, lawnmower path planner, GPS-based target localisation, web-based ground station, and the simulation environment. Also built the model training pipeline and developed the progressive testing framework.

\textbf{\textcolor{red}{Robin [Surname]} --- Flight Dynamics and State Machine Testing.}
MSc student in Aerospace Engineering. Contributed to flight dynamics analysis and testing of the autonomous state machine, verifying state transitions and flight parameter tuning in both simulation and bench testing. Supported integration of the Cube Orange+ flight controller with the mission software and assisted with flight mode configuration (GUIDED, LOITER, STABILIZE, LAND).

\textbf{\textcolor{red}{[TEAMMATE~3]} --- Designated Pilot and RC Systems.}
MSc student in Aerospace Engineering with prior experience in radio-controlled aircraft. Served as the company's designated pilot, responsible for RC transmitter configuration, kill-switch setup, and manual override procedures. Maintained readiness to assume manual control during all autonomous flight tests, ensuring compliance with the requirement that a pilot with override authority be present at all times.

\textbf{\textcolor{red}{[TEAMMATE~4]} --- Hardware Integration.}
MSc student in Mechanical Engineering. Led the physical assembly of the drone platform, including airframe construction, wiring of the Cube Orange+ to the Raspberry Pi (serial TELEM2 connection), camera mounting, GPS antenna placement, and compass calibration. Ensured that all electrical connections were secure and that the payload (Pi, camera, first-aid release mechanism) was mounted within the airframe's centre-of-gravity envelope.

\textbf{\textcolor{red}{[TEAMMATE~5]} --- Project Management and Documentation.}
MSc student in Aerospace Engineering. Coordinated project scheduling, booked flight slots at the Fenswood site, managed risk assessment paperwork, and maintained deliverable tracking against the module timeline. Led the report writing effort, compiled the requirements verification evidence, and ensured that all deliverables (D1--D6) were submitted on time.

\subsection*{Contribution Summary}

Table~\ref{tab:contributions} summarises the key contributions of each team member across the major project areas.

\begin{table}[h]
\centering
\caption{Team member contributions by project area.}
\label{tab:contributions}
\small
\begin{tabular}{p{3.2cm} p{10.5cm}}
\toprule
\textbf{Team Member} & \textbf{Key Contributions} \\
\midrule
Dmytro Chuprynyuk &
Autonomous software stack (state machine, CV pipeline, path planner, target localisation, ground station, simulation environment). Custom YOLOv8n model training and dataset generation. Progressive testing framework. Pi hardware integration (camera, TFLite inference, Cube serial link). Lens calibration and FOV characterisation. \\
\addlinespace
\textcolor{red}{Robin [Surname]} &
Flight dynamics analysis. State machine testing in simulation and on bench. Flight parameter tuning (altitude, speed, descent profile). Cube flight-mode configuration and verification. \\
\addlinespace
\textcolor{red}{[TEAMMATE~3]} &
RC transmitter configuration and binding. Kill-switch and flight-mode switch setup. Manual override procedures. Designated pilot for all flight tests. \\
\addlinespace
\textcolor{red}{[TEAMMATE~4]} &
Airframe assembly and wiring. Cube--Pi serial connection. Camera and GPS antenna mounting. Compass calibration. Payload integration and CG verification. \\
\addlinespace
\textcolor{red}{[TEAMMATE~5]} &
Project scheduling and milestone tracking. Flight-slot booking and risk assessments. Report compilation and editing. Requirements verification evidence. Deliverable management (D1--D6). \\
\bottomrule
\end{tabular}
\end{table}

\subsection*{Approach and Objectives}

Our approach rests on three principles. \textbf{First, simulation before flight}: the full autonomy stack---state machine, computer vision, path planning, geofencing---was developed and tested end-to-end in a Software-In-The-Loop (SITL) environment before any hardware was powered on, ensuring that logic errors were caught at zero hardware risk. \textbf{Second, modular single-interface design}: a single function (\texttt{detect\_in\_image}) encapsulates all computer vision, allowing three model retrains and backend changes with zero modifications to the mission orchestrator. \textbf{Third, progressive testing}: a five-tier framework (simulation $\to$ Docker $\to$ bench $\to$ passive flight $\to$ autonomous flight) gates each integration step, catching defects at the lowest-cost tier possible.

The project objectives, derived from the twelve module requirements, are:
\begin{enumerate}[nosep]
    \item Autonomously search a polygon area with 100\% coverage while avoiding the SSSI no-fly zone.
    \item Detect a casualty surrogate from $\geq$\SI{30}{\metre} altitude using onboard AI at $\geq$\SI{4}{FPS}.
    \item Localise the target to $<$\SI{5}{\metre} GPS accuracy using vision-based estimation.
    \item Land within 5--10\,m of the target, deploy a first-aid payload, and return to the take-off location.
    \item Validate all capabilities through progressive testing with quantitative evidence.
\end{enumerate}

\subsection*{Report Structure}

The report is organised as follows. Section~\ref{sec:rationale} presents the \emph{Design Rationale}, including STEEPLE analysis and four MCDA trade studies that justify the key architectural choices. Sections~\ref{sec:sysdesc}--\ref{sec:groundstation} provide the \emph{System Description}: hardware platform, computer vision pipeline, path planning, state machine, target localisation, and ground station. Section~\ref{sec:testing} covers the progressive testing methodology, simulation environment, and defect discovery results. Section~\ref{sec:requirements} presents \emph{Requirements Verification} with evidence of compliance for all twelve requirements. Section~\ref{sec:evaluation} provides a plus/delta \emph{Evaluation} of technical performance, including comparison with published SAR systems and a critical assessment of design choices. Appendices provide the full state-transition table, detection samples, calibration data, and the complete defect register.
